\section{Conclusion}
Phase-scoped QB expresses one-query attention through paired query/mask views while retaining surrounding batching. Its phase dispatch, head mapping, output recovery and graph-cache conditions integrate the representation with a stateful speculative decoder. The original three-family experiments establish bitwise local correspondence and 2.83--6.17\% lower speculative request latency than PQ. A matched Qwen3 study retains the practical benefit at $k=3$ with both resident and fresh contexts.

The registered three-date study measures the repair-scope trade-off. Qwen3 moves from a draft value of 0.835 under GS to 1.188 under QB, while SmolLM2 remains faster without drafting at 0.805. Both scope interactions are positive on every date mean. Repair scope changes the full-request trade-off, and the resulting draft decision remains model- and workload-dependent.

Reference eligibility also depends on the state supplied to local computation. Two controlled Vulkan cases retain the same natural serial endpoint while a producer-state change switches PQ/QB from reference divergence to preservation. The two fixed-reference cases require a producer--consumer qualification relation that independent component pass labels cannot express. Together with the matched cost studies, they separate reference eligibility, QB's representation benefit and the decision to enable drafting on the tested workloads.
